% ============================================================
% Capítulo 5: Prompts para Arquitetura e Código
% ============================================================
\chapter{Prompts para Arquitetura e Código}

\section{Introdução}

% [Placeholder: A arte de pedir à IA]
% - Prompt engineering como habilidade de engenharia
% - Contexto é tudo
% - Estrutura de um bom prompt

\section{Anatomia de um Bom Prompt}

% [Placeholder: Elementos essenciais]
% - Contexto e papel da IA
% - Descrição clara do problema
% - Restrições e requisitos
% - Formato de saída esperado
% - Exemplos (few-shot prompting)

\section{Prompts para Design Arquitetural}

% [Placeholder: Usando IA para arquitetura]
% - Prompt para explorar alternativas arquiteturais
% - Prompt para analisar trade-offs
% - Prompt para gerar ADRs
% - Prompt para identificar riscos no design
% - Exemplos práticos

\section{Prompts para Geração de Código}

% [Placeholder: Gerando código com qualidade]
% - Especificando interfaces antes da implementação
% - Definindo contratos e tipos
% - Solicitando código com testes
% - Pedindo explicações junto com o código
% - Exemplos práticos

\section{Prompts para Revisão de Código}

% [Placeholder: IA como revisor]
% - Prompt para encontrar bugs
% - Prompt para avaliar segurança
% - Prompt para sugerir melhorias de design
% - Prompt para verificar legibilidade
% - Exemplos práticos

\section{Prompts para Refatoração}

% [Placeholder: Refatorando com assistência]
% - Identificando code smells
% - Solicitando refatorações específicas
% - Mantendo comportamento durante refatoração
% - Exemplos práticos

\section{Padrões Avançados de Prompting}

% [Placeholder: Técnicas avançadas]
% - Chain of Thought prompting
% - Tree of Thoughts
% - Iteração e refinamento
% - Prompt chaining para tarefas complexas
% - System prompts e personalização

\section{Anti-padrões de Prompting}

% [Placeholder: O que não fazer]
% - Prompts vagos e genéricos
% - Assumir que a IA entende contexto implícito
% - Aceitar a primeira resposta sem validação
% - Não fornecer exemplos
% - Ignorar limitações do modelo

\section{Conclusão}

% [Placeholder: Resumo]
% - Prompts são especificações formais
% - Qualidade do prompt = qualidade do resultado

\section*{Exercícios}

% [Placeholder: Exercícios]
% 1. Escreva um prompt para gerar a arquitetura de um sistema de
%    autenticação. Compare com uma arquitetura que você desenharia.
% 2. Crie um prompt para revisar um trecho de código complexo.
%    Avalie a qualidade da análise.
% 3. Experimente chain of thought em um problema de design.
%    Documente o processo.
